\section{Homomorphisms}
\textbf{Conventions.} Permutations compose from right to left. The identity is $e$. A group generated by an element $a$ is denoted $\gen{a}$.

\exercise{13.1}{Inclusion of the integers in the reals}
Yes. Both groups use addition, so the required identity is $\phi(m+n)=\phi(m)+\phi(n)$. The left side is $m+n$ and the right side is also $m+n$. Thus the inclusion preserves the operation and is a homomorphism.

\exercise{13.2}{The greatest-integer map}
No. Choose $x=y=3/5$. Then $\phi(x)=\phi(y)=0$, but $\phi(x+y)=\lfloor6/5\rfloor=1$. Hence $\phi(x+y)\ne\phi(x)+\phi(y)$, so the additive homomorphism law fails.

\exercise{13.3}{Absolute value under multiplication}
Yes. Absolute value sends nonzero reals to nonzero reals, so the map has the stated codomain. For every $x,y\in\R^*$,
\[\phi(xy)=|xy|=|x||y|=\phi(x)\phi(y).\]
The kernel is $\{1,-1\}$, showing that a homomorphism need not be injective.

\exercise{13.4}{Reduction from $\Z_6$ to $\Z_2$}
Yes. Reduction of an integer modulo $2$ is unchanged if a multiple of $6$ is added, because $6$ is even. Thus reduction descends to $\Z_6$, and reducing the sum of two representatives modulo $2$ gives the sum of their reductions. Formally $\phi([a+b]_6)=[a+b]_2=[a]_2+[b]_2$, proving the homomorphism law.

\exercise{13.5}{Remainder parity on $\Z_9$}
No. Using the canonical representatives, take $8$ and $1$. In $\Z_9$ their sum is $0$, so $\phi(8+1)=\phi(0)=0$. But $\phi(8)+\phi(1)=0+1=1$ in $\Z_2$. Subtracting the modulus $9$ changes parity, causing the homomorphism law to fail.

\exercise{13.6}{Exponentiation from addition to multiplication}
Yes. The values $2^x$ are positive and hence nonzero reals. The domain operation is addition and the target operation is multiplication, so the correct identity is
\[\phi(x+y)=2^{x+y}=2^x2^y=\phi(x)\phi(y).\]
The change in operations is part of the homomorphism definition and must be kept explicit.

\exercise{13.7}{A coordinate inclusion into a direct product}
Yes. The map puts $g_i$ in its own coordinate and identity elements in all others. Multiplying $\phi_i(a)$ and $\phi_i(b)$ multiplies only the $i$th coordinates nontrivially, giving $ab$ there. Consequently $\phi_i(a)\phi_i(b)=\phi_i(ab)$. It is also injective because equality of the product vectors forces equality of their $i$th coordinates.

\exercise{13.8}{The inverse map}
For an arbitrary group, the answer depends on commutativity. One has
\[\phi(ab)=(ab)^{-1}=b^{-1}a^{-1},\qquad\phi(a)\phi(b)=a^{-1}b^{-1}.\]
Thus it is a homomorphism if the group is abelian. Conversely, if it is a homomorphism, these expressions are equal for every $a,b$; taking inverses gives $ab=ba$, so the group must be abelian. For example it fails in $S_3$ using two noncommuting transpositions. The exact answer is: if and only if $G$ is abelian.

\exercise{13.9}{Taking a second derivative}
Yes. The second derivative of a function having derivatives of all orders again has derivatives of all orders, so the map stays in $F$. Differentiation is additive, hence
\[\phi(f+g)=(f+g)''=f''+g''=\phi(f)+\phi(g).\]
This is the homomorphism law for the additive function groups.

\exercise{13.10}{Integration as an additive map}
Yes. A continuous function has a finite Riemann integral on $[0,4]$. Linearity of the integral gives
\[\phi(f+g)=\int_0^4(f(x)+g(x))\,dx=\int_0^4f(x)\,dx+\int_0^4g(x)\,dx=\phi(f)+\phi(g).\]
Thus the map preserves the additive operations.

\exercise{13.11}{Multiplying a function by three}
Yes. For pointwise addition,
\[\phi(f+g)=3(f+g)=3f+3g=\phi(f)+\phi(g).\]
The equality holds at every real input, so it is an equality of functions. There is no requirement here that the functions be continuous or differentiable.

\exercise{13.12}{Determinant on an additive matrix group}
For $n\ge2$, no. The domain operation is matrix addition, so a homomorphism would require $\det(A+B)=\det A+\det B$. Taking $A=B=I_n$ gives $\det(2I_n)=2^n$, while $\det I_n+\det I_n=2$; these are unequal for $n\ge2$. The familiar multiplicative determinant identity concerns matrix multiplication and does not apply to this domain operation.
For the exceptional case $n=1$, the determinant is just the single entry, so the map is additive and is a homomorphism. Thus a blanket negative answer implicitly assumes $n>1$.

\exercise{13.13}{Trace on an additive matrix group}
Yes. The diagonal entries of $A+B$ are $a_{ii}+b_{ii}$, so
\[\operatorname{tr}(A+B)=\sum_i(a_{ii}+b_{ii})=\sum_i a_{ii}+\sum_i b_{ii}=\operatorname{tr}(A)+\operatorname{tr}(B).\]
This proves the additive homomorphism law for every positive matrix size.

\exercise{13.14}{Trace on the multiplicative general linear group}
No. A homomorphism from a multiplicative group to additive $\R$ must send the multiplicative identity to $0$. But $\operatorname{tr}(I_n)=n\ne0$ for $n\ge1$. Equivalently, $I_nI_n=I_n$ would require $n=n+n$, which is false. The same formula that works on the additive matrix group fails for this different operation.

\exercise{13.15}{Integration on a multiplicative function group}
The map is well defined: a continuous nonvanishing function on $\R$ has constant sign, so its integral on $[0,1]$ is nonzero. But it is not multiplicative. Take
\[f(x)=g(x)=2+\sin(2\pi x),\]
which is everywhere positive. Then $\phi(f)=\phi(g)=2$, while
\[\phi(fg)=\int_0^1\bigl(4+4\sin(2\pi x)+\sin^2(2\pi x)\bigr)\,dx=4+0+\tfrac12=\tfrac92.\]
This differs from $\phi(f)\phi(g)=4$, so the homomorphism law fails.

\exercise{13.16}{Kernel of the parity homomorphism}
The map of Example 13.3 sends even permutations to $0\in\Z_2$ and odd permutations to $1$. The target identity is $0$, so its inverse image is exactly the even permutations:
\[\Ker\phi=A_3=\{e,(1,2,3),(1,3,2)\}=\{\rho_0,\rho_1,\rho_2\}.\]

\exercise{13.17}{An integer homomorphism into $\Z_7$}
Since $1$ generates $\Z$ and $\phi(1)=4$, additivity forces $\phi(n)=4n\pmod7$ for every integer $n$. Therefore $\phi(25)=100\equiv2\pmod7$. The kernel consists of integers with $7\mid4n$. Since $\gcd(4,7)=1$, this is equivalent to $7\mid n$, so $\Ker\phi=7\Z$.

\exercise{13.18}{An integer homomorphism into $\Z_{10}$}
The image of $n$ is $6n\pmod{10}$. Thus $\phi(18)=108\equiv8\pmod{10}$. Also $10\mid6n$ if and only if $5\mid3n$, which is equivalent to $5\mid n$. Consequently $\Ker\phi=5\Z$.

\exercise{13.19}{An integer homomorphism into $S_8$}
Let $g=(1,4,2,6)(2,5,7)$. Its factors overlap, so first compute their composition:
\[1\mapsto4\mapsto2\mapsto5\mapsto7\mapsto6\mapsto1,
\qquad g=(1,4,2,5,7,6).\]
It has order $6$, and $\phi(n)=g^n$, so $\Ker\phi=6\Z$. Since $20\equiv2\pmod6$,
\[\phi(20)=g^2=(1,2,7)(4,5,6),\]
with the remaining symbols fixed.

\exercise{13.20}{A homomorphism from $\Z_{10}$ to $\Z_{20}$}
The formula is $\phi(n)=8n\pmod{20}$, and it is consistent with the domain modulus because $8\cdot10=80\equiv0\pmod{20}$. The kernel condition $20\mid8n$ reduces to $5\mid n$, so in $\Z_{10}$ the kernel is $\{0,5\}$. Finally $\phi(3)=24\equiv4\pmod{20}$.

\exercise{13.21}{A cyclic homomorphism into $S_8$}
Put $g=(2,5)(1,4,6,7)$. The disjoint cycles have lengths $2,4$, so $g$ has order $4$. Hence $\phi(n)=g^n$ is well defined on $\Z_{24}$, since $4\mid24$. Its kernel consists of multiples of $4$:
\[\Ker\phi=\{0,4,8,12,16,20\}.\]
Because $14\equiv2\pmod4$, squaring the cycles gives
\[\phi(14)=(1,6)(4,7).\]

\exercise{13.22}{A map from two integer coordinates to one}
By additivity,
\[\phi(m,n)=m\phi(1,0)+n\phi(0,1)=3m-5n.\]
Thus $\phi(-3,2)=-9-10=-19$. For the kernel, solve $3m=5n$. Coprimality forces $m=5k$ and then $n=3k$, with $k\in\Z$. Therefore
\[\Ker\phi=\{(5k,3k):k\in\Z\}=\gen{(5,3)}.\]

\exercise{13.23}{A map between integer lattices}
The generator images give
\[\phi(m,n)=m(2,-3)+n(-1,5)=(2m-n,-3m+5n).\]
Therefore $\phi(4,6)=(8-6,-12+30)=(2,18)$. If the image is $(0,0)$, the first coordinate gives $n=2m$, and the second then gives $-3m+10m=7m=0$. Hence $m=n=0$, so the kernel is $\{(0,0)\}$.

\exercise{13.24}{A two-generator homomorphism to $S_{10}$}
Write $a=(3,5)(2,4)$ and $b=(1,7)(6,10,8,9)$. Their supports are disjoint, so they commute, making the formula $\phi(m,n)=a^mb^n$ consistent with the abelian domain. The orders are $2$ and $4$.
If $a^mb^n=e$, restriction to the support of $a$ gives $2\mid m$, and restriction to that of $b$ gives $4\mid n$. Conversely those divisibilities clearly suffice. Thus
\[\Ker\phi=2\Z\times4\Z.\]
For the specified point, $a^3=a$ and $b^{10}=b^2=(6,8)(10,9)$, so
\[\phi(3,10)=(3,5)(2,4)(6,8)(9,10).\]

\exercise{13.25}{Homomorphisms of $\Z$ onto itself}
Every homomorphism has the form $\phi(n)=kn$, where $k=\phi(1)\in\Z$. Its image is $k\Z$. This is all of $\Z$ precisely when $k=1$ or $k=-1$. Thus there are exactly two surjective homomorphisms, $n\mapsto n$ and $n\mapsto-n$.

\exercise{13.26}{Homomorphisms of $\Z$ into itself}
For each integer $k$, the formula $\phi_k(n)=kn$ is additive, so is a homomorphism. Conversely additivity and the generator $1$ force this formula for every homomorphism, including negative inputs by $\phi(-n)=-\phi(n)$. Therefore the homomorphisms are parametrized by $k\in\Z$: countably infinitely many.

\exercise{13.27}{Homomorphisms of $\Z$ into $\Z_2$}
The image of $1$ is either $0$ or $1$, and that choice determines the map on every integer. Both choices work: $n\mapsto0$ and $n\mapsto n\bmod2$ preserve addition. Thus there are exactly two homomorphisms.

\exercise{13.28}{When is left multiplication a homomorphism?}
A homomorphism must send identity to identity. Here $\phi_g(e)=ge=g$, so this requires $g=e$. Conversely, when $g=e$, $\phi_g$ is the identity map and is a homomorphism. Thus precisely the identity element qualifies. Being a permutation of $G$ is much weaker than being a group homomorphism.

\exercise{13.29}{When is conjugation a homomorphism?}
For every $g\in G$. Indeed
\[\phi_g(xy)=gxyg^{-1}=(gxg^{-1})(gyg^{-1})=\phi_g(x)\phi_g(y),\]
where the middle $g^{-1}g$ cancels. Its inverse map is conjugation by $g^{-1}$, so it is even an automorphism for every $g$.

\exercise{13.30}{Definition of a group homomorphism}
A complete definition specifies the algebraic structures and the quantifier: a homomorphism from a group $G$ to a group $G'$ is a function $\phi:G\to G'$ satisfying $\phi(xy)=\phi(x)\phi(y)$ for all $x,y\in G$. The first product is in $G$ and the second is in $G'$. Merely saying ``a map'' without domains, operations, and a universal condition leaves the original statement incomplete.

\exercise{13.31}{Definition of the kernel}
The definition is correct. For a group homomorphism $\phi:G\to G'$ with target identity $e'$, the kernel is the inverse image of that identity:
\[\Ker\phi=\{x\in G:\phi(x)=e'\}.\]
It always contains the identity of $G$ and is a normal subgroup; the identity used in the test is the target identity, which need not have the same notation or value as the domain identity.

\exercise{13.32}{True or false}
\begin{enumerate}[label=(\alph*)]
\item True: $A_n$ is the kernel of the sign homomorphism (with the trivial small cases included).
\item True: the constant map $g\mapsto e'$ is always a homomorphism.
\item False: a nontrivial group mapping constantly to the target identity has many points in its kernel.
\item True: $\phi(x)=\phi(y)$ is equivalent to $xy^{-1}\in\Ker\phi$, giving injectivity exactly for a trivial kernel.
\item False: an image of a six-element group has order dividing $6$, so cannot have order $4$.
\item False: a function on six points has at most six image points, so cannot have an image of size $12$.
\item True: $\Z_6\to\Z_{12}$, $[k]_6\mapsto[2k]_{12}$, is a homomorphism, even injective.
\item True: at least the trivial homomorphism exists; a nontrivial example is $\Z_6\to\Z_{10}$, $[k]\mapsto[5k]$.
\item False: the domain identity always maps to the target identity, so the kernel is never empty.
\item False: $\Z_2\to\R^*$, $[k]\mapsto(-1)^k$, is a nontrivial homomorphism into an infinite group.
\end{enumerate}

\exercise{13.33}{A nontrivial map from $\Z_{12}$ to $\Z_5$?}
No such map exists. Its image order must divide both the domain order $12$ and the target order $5$. Their gcd is $1$, so the image has order $1$ and the map is trivial. Equivalently, if $a=\phi(1)$, then $12a=0$ in $\Z_5$ forces $a=0$.

\exercise{13.34}{A nontrivial map from $\Z_{12}$ to $\Z_4$}
Use reduction modulo $4$: $\phi([k]_{12})=[k]_4$. It is well defined because changing $k$ by $12$ does not change its residue modulo $4$, and it preserves addition. It is nontrivial because $\phi(1)=1\ne0$, and is in fact surjective.

\exercise{13.35}{A nontrivial map between two direct products}
For the source $\Z_2\times\Z_4$ and target $\Z_2\times\Z_5$, define $\phi(a,b)=(a,0)$. Addition is coordinatewise, so
\[\phi((a,b)+(c,d))=(a+c,0)=(a,0)+(c,0).\]
Thus it is a homomorphism. The value $\phi(1,0)=(1,0)$ is nonidentity, so it is nontrivial.

\exercise{13.36}{A nontrivial map from $\Z_3$ to $\Z$?}
Impossible. If $k=\phi(1)$, then $3k=\phi(3\cdot1)=\phi(0)=0$ in the integers. This forces $k=0$, so every element maps to zero. The obstruction is that a finite-order element cannot have a nonzero image in the torsion-free group $\Z$.

\exercise{13.37}{A nontrivial map from $\Z_3$ to $S_3$}
Let $r=(1,2,3)$ and define $\phi([k]_3)=r^k$. Since $r^3=e$, this is independent of the representative. The exponent law gives $\phi(a+b)=r^{a+b}=r^ar^b=\phi(a)\phi(b)$. Its image is $A_3$, and $\phi(1)=r\ne e$, so it is nontrivial.

\exercise{13.38}{A nontrivial map from $\Z$ to $S_3$}
Take $\phi(k)=(1,2)^k$. Integer powers satisfy $\phi(k+\ell)=\phi(k)\phi(\ell)$, including negative exponents. The map is nontrivial because $\phi(1)=(1,2)\ne e$. Its image is the order-$2$ subgroup $\{e,(1,2)\}$.

\exercise{13.39}{A nontrivial map from $\Z^2$ to $2\Z$}
Define $\phi(m,n)=2m$. Its values lie in the even integers, and
\[\phi((m,n)+(r,s))=2(m+r)=2m+2r=\phi(m,n)+\phi(r,s).\]
It is nontrivial since $\phi(1,0)=2$, and is even onto $2\Z$.

\exercise{13.40}{A nontrivial map from $2\Z$ to $\Z^2$}
For an even integer $2k$, define $\phi(2k)=(k,0)$. The integer $k$ is unique, so the map is well defined. Then $\phi(2k+2\ell)=(k+\ell,0)=(k,0)+(\ell,0)$. It is nontrivial because $\phi(2)=(1,0)$.

\exercise{13.41}{A nontrivial map from $D_4$ to $S_3$}
View $D_4$ as permutations of its four vertices. Define $\phi(g)=e$ when $g$ is even and $\phi(g)=(1,2)$ when $g$ is odd. Parity adds modulo $2$, and $(1,2)^2=e$, so this is a homomorphism into $S_3$. A quarter-turn is a $4$-cycle and is odd, so it has nonidentity image, proving the map is nontrivial.

\exercise{13.42}{A nontrivial map from $S_3$ to $S_4$}
Extend a permutation of $\{1,2,3\}$ by fixing $4$. The extension of a composite equals the composite of the extensions, because this is true on $1,2,3$ and all maps fix $4$. Thus the extension is an injective homomorphism and is nontrivial.

\exercise{13.43}{A nontrivial map from $S_4$ to $S_3$}
Send each even permutation to $e\in S_3$ and each odd permutation to $(1,2)$. This is the sign homomorphism with its two values relabeled by the order-$2$ subgroup $\{e,(1,2)\}$ of $S_3$. Products preserve parity, so it preserves the group operation. It is nontrivial because a transposition has image $(1,2)$. Surjectivity onto $S_3$ is not required by the exercise.

\exercise{13.44}{Image order divides a finite domain order}
Let $K=\Ker\phi$. For an image element $\phi(a)$, its full fiber is $aK$: indeed $\phi(x)=\phi(a)$ if and only if $\phi(a^{-1}x)=e'$, or $a^{-1}x\in K$. Each fiber therefore has $|K|$ elements by $k\mapsto ak$, and the fibers partition finite $G$. Thus
\[|G|=|\phi[G]|\,|K|.\]
The image is finite and its order divides $|G|$. This proof obtains the count directly without presupposing the isomorphism theorem of the next section.

\exercise{13.45}{Image order divides a finite target order}
The image $\phi[G]$ is a subgroup of $G'$: it contains $e'=\phi(e)$, products satisfy $\phi(a)\phi(b)=\phi(ab)$, and inverses satisfy $\phi(a)^{-1}=\phi(a^{-1})$. Since $G'$ is finite, its subgroup is finite and Lagrange gives $|\phi[G]|\mid|G'|$. The domain need not be finite for this argument.

\exercise{13.46}{Agreement on generators determines a homomorphism}
Every element of the group generated by the $a_i$ is a finite word $a_{i_1}^{\varepsilon_1}\cdots a_{i_k}^{\varepsilon_k}$ with $\varepsilon_j=\pm1$. A homomorphism carries this word to the product of the corresponding image powers, because it preserves products and inverses. Since $\phi(a_i)=\mu(a_i)$ for every generator, the two maps give equal products on every such word. They also agree at the identity. Thus they agree on every element of $G$, meaning $\phi=\mu$. No assumption that the generators are independent is needed.

\exercise{13.47}{A homomorphism with prime-order domain}
If $|G|=p$ is prime, its kernel is a subgroup and must have order $1$ or $p$ by Lagrange. In the first case the kernel is trivial, so the homomorphism is injective: equal images imply the quotient of the inputs lies in that kernel. In the second case the kernel is all of $G$, so every input maps to the target identity and the homomorphism is trivial. These are the only possibilities.

\exercise{13.48}{Kernel of the sign map}
By definition $\operatorname{sgn}(\sigma)=+1$ exactly for even permutations. The identity of the multiplicative target $\{1,-1\}$ is $+1$, so
\[\Ker(\operatorname{sgn}_n)=A_n.\]
For $n\ge2$, a transposition has sign $-1$, so the map is onto both target values. For $n=1$ it is still a homomorphism with kernel $S_1=A_1$, but it is not onto the two-element target. This map is the general version of the parity map $S_3\to\Z_2$.

\exercise{13.49}{Composition of homomorphisms}
For $a,b\in G$, apply first the homomorphism law for $\phi$, then the one for $\psi$:
\[(\psi\phi)(ab)=\psi(\phi(ab))=\psi(\phi(a)\phi(b))=\psi(\phi(a))\psi(\phi(b)).\]
The result is $(\psi\phi)(a)(\psi\phi)(b)$, which is the required law for the composite $G\to G''$.

\exercise{13.50}{An abelian image and commutators}
For any $x,y\in G$,
\[\phi(xyx^{-1}y^{-1})=\phi(x)\phi(y)\phi(x)^{-1}\phi(y)^{-1}.\]
If the image is abelian, the right side cancels to the identity, so every displayed commutator lies in the kernel. Conversely, if it is always in the kernel, the right side equals the identity. Multiplying that equality on the right by $\phi(y)\phi(x)$ gives $\phi(x)\phi(y)=\phi(y)\phi(x)$. Every image element is some $\phi(x)$, so the image is abelian.

\exercise{13.51}{The universal power map from the integers}
For all integers $m,n$, $a^{m+n}=a^ma^n$, so $\phi(n)=a^n$ is a homomorphism. Its image is exactly $\gen a$. If $a$ has infinite order, only the exponent $0$ gives $e$, so the kernel is $\{0\}$. If $a$ has finite order $r$, then $a^n=e$ exactly when $r\mid n$, so the kernel is $r\Z$. This includes $r=1$, when $a=e$ and the kernel is all of $\Z$.

\exercise{13.52}{A fiber is a kernel coset}
Let $H=\Ker\phi$. For $x\in G$,
\[\phi(x)=\phi(a)\iff\phi(xa^{-1})=e'\iff xa^{-1}\in H\iff x\in Ha.\]
The first equivalence uses preservation of inverses and cancellation in the target; the last is the definition of a right coset. Thus the two described sets are equal. Since a kernel is normal, this fiber is also $aH$.

\exercise{13.53}{A two-generator power formula}
The necessary and sufficient condition is $hk=kh$.
If $\phi(m,n)=h^mk^n$ is a homomorphism, the commuting domain elements $(1,0)$ and $(0,1)$ must have commuting images. Their images are $h,k$, proving necessity.
Conversely suppose $h,k$ commute. Their integer powers then commute, including negative powers, and for all integers $m,n,r,s$,
\[\phi(m,n)\phi(r,s)=h^mk^nh^rk^s=h^{m+r}k^{n+s}=\phi(m+r,n+s).\]
This is the additive-domain homomorphism law, proving sufficiency.

\exercise{13.54}{When all two-generator power formulas work}
The condition is that $G$ be abelian. If every pair $h,k$ gives a homomorphism in Exercise 53, its necessary condition says every pair commutes, so $G$ is abelian. Conversely, if $G$ is abelian, every pair commutes and the sufficient calculation in that exercise proves every such formula is a homomorphism.

\exercise{13.55}{Power maps from a finite cyclic group}
The exact condition is
\[\boxed{h^n=e,}\]
equivalently, $h$ has finite order dividing $n$. Necessity follows because $n[1]=[0]$ in $\Z_n$, so a homomorphism must satisfy $h^n=\phi(n[1])=\phi([0])=e$. For sufficiency, if integers $i,j$ differ by a multiple of $n$, then $h^i=h^j$ because $h^n=e$. Thus $\phi([i])=h^i$ is well defined on residue classes, and
\[\phi([i]+[j])=h^{i+j}=h^ih^j=\phi([i])\phi([j]).\]
The source writes $0\le i\le n$, duplicating $[0]=[n]$; the same condition resolves that endpoint consistency. It requires $h^n=e$, not $h=e$: for $n=2$ and $G=\Z_2$, sending $1$ to the nonzero element gives the identity homomorphism and has $h\ne e$.
